\scanpage{186}
On va plutôt essayer de calculer
\[
 \tau'-\tau=\operatorname{Tr}xyz-\operatorname{Tr}yxz=\operatorname{Tr}([x,y]z)=\operatorname{Tr}([y,z]x)=\operatorname{Tr}([z,x].y)
\]

C'est une expression \inserted{multilin.} alternée en $x,y,z$, donc on doit avoir
\[
 \tau'-\unsure{\tau|z}=\deleted{\unsure{\varepsilon|1/2|\operatorname{Tr}}}(a\wedge x\wedge y\wedge z)\otimes e^{-1}
\]
pour $a\in\unsure{A|\mathcal A}$ bien déterminé. Comme cette expression est invariante par $G=\underline{\mathrm{Aut}}\,A$, \deleted{on voit que} et que (si $g\in G(-)$)
\[
 g(a)\wedge g(x)\wedge g(y)\wedge g(z)=a\wedge x\wedge y\wedge z\qquad(=a\wedge g(x)\wedge g(y)\wedge g(z))
\]
on en conclut que $a=ga$, i.e. $a$ inv. sous $G$ i.e. $a=\varepsilon.1$. Donc on doit avoir
\[
 \tau'-\tau=\varepsilon(1\wedge x\wedge y\wedge z)\otimes e^{-1}
\]
et pour calculer $\varepsilon$, il suffit de prendre $x,y,z$ tels que $1\wedge x\wedge y\wedge z$ inv. i.e. $1,x,y,z$ une base \inserted{de $A$}. On prendra

\[
 x_0=\begin{pmatrix}0&0\\1&0\end{pmatrix},\ y_0=\begin{pmatrix}0&1\\0&0\end{pmatrix},\ z_0=[x_0,y_0]=\begin{pmatrix}-1&0\\0&1\end{pmatrix}=1-2\alpha
 \quad\text{où }\alpha=\begin{pmatrix}1&0\\0&0\end{pmatrix}
\]

(NB \(\displaystyle x_0y_0=\begin{pmatrix}0&0\\0&1\end{pmatrix},\ y_0x_0=\begin{pmatrix}1&0\\0&0\end{pmatrix}\))
\deleted{$x_0y_0z_0=\begin{pmatrix}0&0\\0&1\end{pmatrix}$)}

On trouve
\[
 1\wedge x_0\wedge y_0\wedge z_0=-2\,1\wedge x_0\wedge y_0\wedge\alpha
 =-2\alpha'\wedge x_0\wedge y_0\wedge\alpha
 =2\deleted{x\wedge y\wedge\alpha\wedge\alpha'}\inserted{\alpha\wedge\alpha'\wedge x\wedge y}
\]
où
\[
 \alpha=\begin{pmatrix}1&0\\0&0\end{pmatrix},\quad \alpha'=\begin{pmatrix}0&0\\0&1\end{pmatrix}
\]
et \deleted{$x,y$,} $\alpha,\alpha',x_0,y_0$ font base can. de $M_2(\mathbb Z)$. On
% archive sheet number: 185
